\section{Conclusion}
\label{sec:conclusion}

The Mean-Median Difference provides partisan fairness evidence that is simpler to
compute, more stable across election cycles, and more resistant to the geographic
sorting defense than the Efficiency Gap. Its central limitation---no adopted legal
threshold---is reframed as an advantage when the evidence is presented comparatively
rather than as a threshold test: the $\Delta\text{MM}$ between enacted and
\textsc{bisect} maps provides the court with a precise measurement of the mapmaker
contribution that does not require accepting any academic proposal about what
constitutes unconstitutional excess.

The core empirical findings are robust. \textsc{Bisect} reduces NC MM by 70\% (from
0.07 to 0.02), and the pattern holds across WI, TX, and FL with 64--70\% reductions.
The residual in \textsc{bisect} maps ($\approx 0.02$) is the irreducible geographic
sorting baseline attributable to the concentration of Democratic voters in NC's
Triangle and Charlotte urban cores. The enacted-bisect gap of $\approx 0.05$ is
the mapmaker contribution---an increase in vote-share skewness that cannot be
explained by where voters live and must therefore be explained by deliberate
redistricting choices.

MM's invariance to uniform partisan swings is its most important practical advantage.
Expert witnesses facing cross-examination on election-cycle sensitivity can demonstrate
that MM is mathematically immune to the wave-election critique: uniform swings shift
mean and median equally, leaving MM unchanged. This invariance reflects that MM
measures the map's structural partisan composition---the distribution of competitive
and safe districts---rather than the particular partisan outcome of any single election.

For state court expert witnesses, the recommended presentation is straightforward:
compute \textsc{bisect} reference MM, compute enacted MM, report the gap as the
mapmaker contribution, and explain that the gap is consistent across all four election
cycles and invariant to uniform partisan swings. This presentation directly addresses
the geographic sorting defense, provides a non-threshold-dependent framing that courts
can evaluate under their applicable state constitutional standard, and is robust
to the standard cross-examination challenges.

\bigskip

\noindent\textit{The author thanks the redistricting research team for data
access and algorithmic development. All errors are the author's own.}
